% =====================================================================
%  Catatan Presentasi - Abd. Rouf
%  Render: xelatex catatan-presentasi.tex   (jalankan 2x)
% =====================================================================
% !TEX program = xelatex
\documentclass[11pt,a4paper]{article}
\usepackage{fontspec}
\setmainfont{Calibri}
\usepackage[margin=2.2cm]{geometry}
\usepackage{array,longtable,booktabs}
\usepackage[table]{xcolor}
\usepackage{enumitem}
\usepackage{titlesec}
\usepackage[hidelinks]{hyperref}

\definecolor{head}{HTML}{1F3A5F}
\definecolor{band}{HTML}{E8EEF5}
\renewcommand{\arraystretch}{1.35}
\setlength{\parindent}{0pt}
\setlength{\parskip}{5pt}
\titleformat{\section}{\large\bfseries\color{head}}{\thesection.}{6pt}{}
\titleformat{\subsection}{\normalsize\bfseries}{}{0pt}{}
\pagestyle{plain}
\setlist{nosep,leftmargin=*,topsep=2pt}
\newcommand{\term}[1]{\cellcolor{band}\textbf{#1}}
\newcolumntype{K}{>{\raggedright\arraybackslash}p{3.8cm}}
\newcolumntype{V}{>{\raggedright\arraybackslash}p{12.2cm}}
\newcommand{\kepala}[2]{\rowcolor{head}\textcolor{white}{\textbf{#1}} & \textcolor{white}{\textbf{#2}}\\ \midrule}

\begin{document}

\begin{center}
{\Large\bfseries\color{head} CATATAN PRESENTASI ANATOMI PENELITIAN}\\[5pt]
{\bfseries Analisis QoS Jaringan Berbasis \textit{Crowdsourcing} pada Aplikasi Presensi Menggunakan K-Means}\\[6pt]
Abd. Rouf \quad|\quad NIM 223140118
\end{center}

% ------------------------------------------------------------------
\section{Struktur Penelitian (Sesuai Arahan Dosen)}
\begin{tabular}{@{}lp{11.5cm}@{}}
\textbf{Bidang / topik} & Jaringan: kualitas layanan (QoS) jaringan Wi-Fi dan seluler\\
\textbf{Teknik / algoritma} & \textit{Clustering} K-Means (+ standar penilaian QoS TIPHON)\\
\textbf{Penerapan / kasus} & Aplikasi presensi PT Semeru Agung Mandiri sebagai alat ukur\\
\textbf{Metode penelitian} & \textit{Research and Development} (R\&D): membuat modul telemetri, menguji, lalu memakainya\\
\textbf{Instrumen} & Modul telemetri di dalam aplikasi presensi\\
\textbf{Teknik pengumpulan data} & \textit{Crowdsourcing} (\textit{opportunistic sensing}) dari HP banyak karyawan\\
\end{tabular}

Polanya sama dengan contoh dosen ``IoT + Fuzzy'': di sini \textbf{Jaringan + K-Means}. R\&D adalah cara kerjanya, bukan topik keilmuannya.

% ------------------------------------------------------------------
\section{Posisi Crowdsourcing, Telemetri, dan K-Means}
Ketiga istilah ini sering tertukar. Masing-masing punya peran berbeda:
\begin{longtable}{KV}
\toprule
\kepala{Unsur}{Perannya di penelitian (dan letaknya di Bab 3)}
\endhead
\term{R\&D} & \textbf{Metode penelitian}, yaitu kerangka besar cara kerja: analisis, rancang, implementasi, validasi, pengumpulan \& analisis. (Subbab Metode Penelitian) \\
\term{Telemetri} & \textbf{Instrumen/alat ukur}: modul di aplikasi yang mengukur kualitas jaringan lalu mengirim hasilnya ke server. Bekerja di \textbf{satu HP}. (Subbab Instrumen Penelitian) \\
\term{Crowdsourcing} & \textbf{Teknik pengumpulan data}: telemetri dijalankan di \textbf{HP banyak karyawan sekaligus}, lalu datanya digabung. Menjawab ``datanya dari mana?''. (Subbab Teknik Pengumpulan Data) \\
\term{TIPHON} & \textbf{Standar penilaian}: memberi kategori per metrik (sangat bagus s.d. buruk). (Subbab Teknik Analisis Data) \\
\term{K-Means} & \textbf{Metode analisis data}: mengelompokkan lokasi $\times$ jaringan $\times$ jam menjadi klaster baik/sedang/buruk. Menjawab ``lokasi mana yang buruk?''. (Subbab Teknik Analisis Data) \\
\bottomrule
\end{longtable}
\textbf{Alur:} telemetri mengukur $\rightarrow$ \textit{crowdsourcing} mengumpulkan dari banyak HP $\rightarrow$ TIPHON memberi kategori $\rightarrow$ K-Means mengelompokkan $\rightarrow$ peta kualitas jaringan per lokasi.

\textbf{Ibaratnya:} telemetri adalah termometer yang otomatis melapor; \textit{crowdsourcing} adalah membagikan termometer itu ke 90 orang; K-Means adalah mengelompokkan hasilnya menjadi ``panas'', ``hangat'', dan ``dingin''.

\textbf{Kalau ditanya ``metodenya apa?''}, jawab berurutan: \textit{metode penelitian R\&D, teknik pengumpulan data crowdsourcing, metode analisis K-Means.}

% ------------------------------------------------------------------
\section{Protokol Jaringan yang Terlibat}
Penelitian jaringan yang ``dibungkus'' ke aplikasi biasanya ditanya: \textbf{protokol apa yang dipakai/diukur, dan di lapisan mana?}
\begin{longtable}{KV}
\toprule
\kepala{Lapisan}{Protokol di aplikasi presensi}
\endhead
\term{Aplikasi} & \textbf{HTTPS} (HTTP/1.1 di atas TLS) untuk REST API PocketBase, termasuk \textit{probe} \texttt{GET /api/health} dan unggah data presensi. \textbf{SSE} (\textit{Server-Sent Events}) untuk fitur \textit{realtime} PocketBase. DNS untuk mencari alamat server. FCM untuk notifikasi. \\
\term{Transport} & \textbf{TCP} dengan TLS 1.2/1.3 (di bawah HTTPS dan SSE). \\
\term{Jaringan} & IPv4/IPv6. \\
\term{Data link / fisik} & \textbf{IEEE 802.11} (Wi-Fi) dan seluler (4G LTE/5G). Jenis jaringan dicatat dari \texttt{connectivity\_plus}. \\
\bottomrule
\end{longtable}
\textbf{Yang perlu dikatakan terus terang:}
\begin{itemize}
\item \textit{Delay} diukur sebagai \textbf{RTT tingkat aplikasi (HTTP)}, bukan ping ICMP, karena aplikasi Android biasa tidak bisa mengirim paket ICMP mentah.
\item \textit{Packet loss} diganti \textbf{tingkat gagal} permintaan HTTP sebagai \textit{proxy}.
\item Validasi dengan PCAPdroid mengukur RTT di tingkat TCP sebagai pembanding.
\item Pengukuran bersifat \textbf{hibrida}: pasif dari komunikasi presensi yang sudah berjalan, ditambah \textit{probe} aktif yang sangat ringan (3 permintaan, paling sering tiap 15 menit).
\end{itemize}
\textbf{Rujukan tambahan untuk Bab 2 (opsional):} RFC 2330 (kerangka IPPM), RFC 2681 (\textit{round-trip delay}), RFC 3393 (variasi \textit{delay}/\textit{jitter}), RFC 3550 (rumus \textit{jitter} RTP), ITU-T Y.1540/Y.1541 (parameter dan kelas QoS IP), dan Howe (2006) untuk asal istilah \textit{crowdsourcing}.

% ------------------------------------------------------------------
\section{Penjelasan Singkat (30 Detik)}
\fbox{\parbox{\dimexpr\linewidth-2\fboxsep-2\fboxrule}{\itshape
``Karyawan PT Semeru Agung Mandiri absen lewat aplikasi dari kantor, posko, dan gudang yang sinyalnya berbeda-beda, tetapi perusahaan tidak punya data lokasi mana yang jaringannya buruk. Saya menambahkan modul kecil di aplikasi presensi yang \textbf{diam-diam dan anonim} mencatat kualitas jaringan setiap kali aplikasi menghubungi server, tanpa mengubah cara karyawan absen. Hasilnya dinilai dengan standar QoS TIPHON, lalu dikelompokkan dengan \textbf{K-Means} menjadi kategori baik, sedang, dan buruk per lokasi dan per jam. Jadi aplikasi presensi sekaligus menjadi alat ukur jaringan.''}}

% ------------------------------------------------------------------
\section{Manfaat Penelitian}
\subsection{Bagi perusahaan}
\begin{itemize}
\item \textbf{Peta kualitas jaringan per lokasi kerja} tanpa harus mengirim orang mengukur ke setiap posko.
\item \textbf{Dasar keputusan investasi}: tahu posko mana yang perlu dipasang Wi-Fi, \textit{repeater}, atau operator seluler lain, sehingga biaya tepat sasaran.
\item \textbf{Penjelasan atas keluhan presensi}: kalau karyawan bilang ``absen tidak terkirim'', HRD bisa mengecek apakah lokasi dan jam itu memang bermasalah.
\item \textbf{Pemantauan berkelanjutan}: modul tetap bisa dinyalakan setelah penelitian selesai untuk memantau perubahan kualitas jaringan.
\item \textbf{Bahan perbaikan aplikasi}: tahu berapa lama dan seberapa sering HP offline, sehingga pengaturan sinkronisasi dan ukuran foto bisa disesuaikan dengan data, bukan perkiraan.
\end{itemize}
\subsection{Bagi karyawan}
\begin{itemize}
\item Lokasi bersinyal buruk bisa diperbaiki, sehingga absen lebih lancar.
\item Tidak ada beban tambahan: tidak perlu mengisi apa pun dan data dibuat anonim.
\end{itemize}
\subsection{Bagi peneliti selanjutnya}
\begin{itemize}
\item \textbf{Contoh metode baru}: aplikasi operasional apa pun (kasir, kurir, presensi) bisa dijadikan alat ukur QoS dengan cara yang sama.
\item \textbf{Desain privasi yang bisa ditiru}: anonim, ringan, dan bisa dimatikan dari server.
\item \textbf{Bahan pengembangan}: data bisa dilanjutkan dengan prediksi kualitas jaringan (misalnya \textit{machine learning}), penjadwalan sinkronisasi pintar, atau perbandingan operator seluler.
\item \textbf{Perbandingan metode}: hasil K-Means dari data aplikasi bisa dibandingkan dengan pengukuran manual Wireshark pada penelitian QoS lain.
\end{itemize}
\subsection{Bagi keilmuan}
Menambah bukti bahwa pengukuran QoS berbasis \textit{crowdsourcing} pasif dapat diterapkan pada skala kecil (satu perusahaan) dan dikombinasikan dengan K-Means.

% ------------------------------------------------------------------
\section{Apakah Penelitian Ini Pernah Ada?}
\begin{longtable}{KV}
\toprule
\kepala{Yang sudah ada}{Contoh}
\endhead
\term{QoS + TIPHON} & Sangat banyak di jurnal Indonesia, tetapi diukur manual dengan Wireshark di satu tempat (Fiqria 2024, Rosid 2023, Ramayanti 2024, Basmar 2026). \\
\term{K-Means + kualitas jaringan} & Ada, tetapi memakai data pihak ketiga (Ookla) atau data perangkat jaringan (Irwansyah 2025, Mubarak 2023). \\
\term{Crowdsourcing / pengukuran pasif} & Ada di luar negeri dengan aplikasi khusus pengukur dan skala jutaan pengguna (Wamser 2021, Wu 2017). \\
\bottomrule
\end{longtable}
\textbf{Yang baru dari penelitian ini:} menggabungkan ketiganya, yaitu mengukur QoS secara pasif dari \textbf{aplikasi operasional perusahaan (presensi)}, lalu mengelompokkannya dengan K-Means per lokasi kerja. Dari pencarian yang sudah dilakukan, kombinasi ini belum ditemukan. Kalimat yang aman untuk dosen: \textit{``Sepanjang penelusuran saya, belum ditemukan penelitian yang memakai aplikasi presensi sebagai alat ukur QoS crowdsourcing.''}

% ------------------------------------------------------------------
\section{Tentang Judul}
\subsection{Judul yang dipakai (versi padat, sesuai permintaan dosen)}
\fbox{\parbox{\dimexpr\linewidth-2\fboxsep-2\fboxrule}{\bfseries
Analisis QoS Jaringan Berbasis \textit{Crowdsourcing} pada Aplikasi Presensi Menggunakan K-Means}}

Judul awal yang disetujui (19 kata): \textit{Analisis Quality of Service Jaringan Berbasis Crowdsourcing pada Aplikasi Presensi Mobile Menggunakan Clustering K-Means di PT Semeru Agung Mandiri}. Versi padat tinggal \textbf{10 kata}, dan \textbf{isi penelitiannya tidak berubah}.

\subsection{Apa yang dipangkas dan kenapa}
\begin{longtable}{KV}
\toprule
\kepala{Dipangkas}{Alasan}
\endhead
\term{Quality of Service $\rightarrow$ QoS} & Singkatan baku di bidang jaringan. Kepanjangannya ditulis di abstrak dan kata kunci. \\
\term{Clustering K-Means $\rightarrow$ K-Means} & K-Means memang algoritma \textit{clustering}, jadi kata ``Clustering'' mengulang. \\
\term{Presensi Mobile $\rightarrow$ Presensi} & ``Mobile'' sudah dijelaskan di Bab 1. \\
\term{di PT Semeru Agung Mandiri} & Dipindah ke Bab 1 sebagai lokasi/objek penelitian. Pedoman TA FTI ITB WIGA tidak mewajibkan nama perusahaan di judul. \\
\bottomrule
\end{longtable}

\subsection{Kenapa judulnya seperti ini}
\begin{longtable}{KV}
\toprule
\kepala{Bagian judul}{Fungsinya}
\endhead
\term{Analisis QoS Jaringan} & \textbf{Tujuan} dan bidang keilmuan: jaringan. Kata ``analisis'' merujuk pada tahap akhir penelitian, yaitu analisis data hasil pemakaian modul telemetri. \\
\term{Berbasis Crowdsourcing} & \textbf{Teknik pengumpulan data} sekaligus kebaruan: data dari banyak HP karyawan, bukan diukur manual di satu titik. \\
\term{pada Aplikasi Presensi} & \textbf{Objek} dan tempat instrumen: aplikasi presensi yang sudah dipakai sehari-hari. \\
\term{Menggunakan K-Means} & \textbf{Metode analisis}, sesuai pola dosen ``IoT + Fuzzy''. \\
\bottomrule
\end{longtable}

\subsection{Variasi kalau ada aturan atau permintaan lain}
\begin{enumerate}
\item \textbf{Nama perusahaan wajib ada:} Analisis QoS Jaringan Berbasis \textit{Crowdsourcing} pada Aplikasi Presensi Menggunakan K-Means di PT Semeru Agung Mandiri.
\item \textbf{Singkatan dilarang:} Analisis Kualitas Layanan Jaringan Berbasis \textit{Crowdsourcing} pada Aplikasi Presensi Menggunakan K-Means.
\item \textbf{Dosen ingin judul terlihat R\&D} (tanyakan dulu, karena fokusnya bergeser ke produk): Pengembangan Modul Telemetri QoS Jaringan Berbasis \textit{Crowdsourcing} pada Aplikasi Presensi Menggunakan K-Means.
\item \textbf{Paling pendek:} Pemetaan Kualitas Jaringan Lokasi Presensi Menggunakan K-Means.
\end{enumerate}
Model ``(Studi Kasus: PT Semeru Agung Mandiri)'' juga boleh kalau dosen lebih suka.

\subsection{Tips agar judul di-ACC}
\begin{itemize}
\item Bawa \textbf{judul padat + variasi} (di atas) supaya bisa langsung menyesuaikan permintaan dosen.
\item Sampaikan pola ``bidang + teknik + metode'' sesuai arahan dosen (bagian 1).
\item Tekankan bahwa datanya pasti bisa diperoleh karena aplikasinya milik sendiri. Dosen biasanya khawatir mahasiswa tidak bisa mendapat data.
\item Istilah asing (\textit{Crowdsourcing}, \textit{Quality of Service}, \textit{Mobile}) ditulis miring. Singkatan QoS dan nama algoritma K-Means tidak perlu miring.
\end{itemize}

% ------------------------------------------------------------------
\section{Hubungan dengan Offline-First}
Penelitian ini \textbf{masih berhubungan} dengan \textit{offline-first}, tetapi \textit{offline-first} sekarang menjadi \textbf{konteks}, bukan fokus.
\begin{itemize}
\item \textbf{Konteks:} aplikasi presensi dibuat \textit{offline-first} justru karena jaringan di lokasi kerja tidak bisa diandalkan. Penelitian ini \textbf{mengukur seberapa buruk} jaringan tersebut.
\item \textbf{Sumber data:} sebagian data telemetri berasal dari percobaan kirim antrean sinkronisasi \textit{offline-first} (berhasil/gagal, lama kirim, ukuran foto).
\item \textbf{Data khusus offline:} telemetri mencatat \textbf{lama HP putus koneksi} (\textit{outage}). Angka ini menunjukkan seberapa sering dan seberapa lama mode \textit{offline} benar-benar dibutuhkan karyawan.
\item \textbf{Manfaat balik:} hasilnya bisa dipakai untuk menyetel sinkronisasi \textit{offline-first} (jeda percobaan ulang, ukuran foto) berdasarkan data nyata.
\end{itemize}
Karena bukan fokus, kata \textit{offline-first} \textbf{tidak perlu} masuk judul. Cukup dibahas di Bab 1 (latar belakang) dan Bab 2 (teori aplikasi yang diteliti).

% ------------------------------------------------------------------
\section{Yang Perlu Disiapkan dan Urutan Kerja}
\begin{enumerate}
\item \textbf{Persetujuan dosen} atas judul dan metode.
\item \textbf{Izin tertulis perusahaan} untuk pengumpulan data kualitas jaringan anonim (wajib, karena data berasal dari HP karyawan).
\item \textbf{Perbaikan keamanan aplikasi dulu.} Audit menemukan APK yang sekarang masih menyimpan kredensial server. Update aplikasi yang membawa telemetri baru boleh dirilis \textbf{setelah} masalah ini diperbaiki (rencana induk fase F1).
\item \textbf{Membuat modul telemetri} di server uji dan aplikasi versi uji.
\item \textbf{Validasi di meja} dengan PCAPdroid.
\item \textbf{Rilis bertahap}: telemetri mati dulu, lalu dinyalakan dari server untuk beberapa orang, kemudian semua.
\item \textbf{Pengumpulan data 4 minggu}, lalu \textbf{analisis} TIPHON + K-Means di laptop.
\end{enumerate}
Rincian teknis setiap langkah ada di \texttt{docs/roadmap/18\_MASTER\_PLAN\_HARDENING\_ARSITEKTUR\_RISET.md} bagian 14 (Track R).

% ------------------------------------------------------------------
\section{Alur Presentasi (Urutan Bicara)}
\begin{enumerate}
\item \textbf{Masalah:} sinyal di lokasi kerja tidak merata, tetapi tidak ada datanya.
\item \textbf{Yang sudah ada:} penelitian QoS biasanya manual di satu titik; penelitian K-Means memakai data pihak ketiga.
\item \textbf{Celah (GAP):} belum ada yang memakai aplikasi operasional perusahaan sebagai alat ukur QoS.
\item \textbf{Solusi:} modul telemetri pasif di aplikasi presensi (anonim, ringan, bisa dimatikan dari server).
\item \textbf{Cara mengukur:} \textit{delay}, \textit{jitter}, tingkat gagal, \textit{throughput} → kategori TIPHON → K-Means.
\item \textbf{Validasi:} hasil telemetri dicocokkan dengan PCAPdroid di meja sebelum dipakai di produksi.
\item \textbf{Manfaat:} peta kualitas jaringan per lokasi sebagai dasar perbaikan jaringan perusahaan.
\end{enumerate}

% ------------------------------------------------------------------
\section{Kamus Istilah (Bahasa Sederhana)}
\begin{longtable}{KV}
\toprule
\kepala{Istilah}{Artinya}
\endhead
\multicolumn{2}{l}{\textbf{A. Istilah jaringan}}\\
\term{QoS (Quality of Service)} & Ukuran \textbf{seberapa bagus kualitas jaringan}, dinilai dari beberapa angka: \textit{delay}, \textit{jitter}, \textit{packet loss}, dan \textit{throughput}. \\
\term{Delay (latensi)} & \textbf{Waktu tunggu} dari kirim sampai dapat balasan. Makin kecil makin bagus. Contoh: HP ``memanggil'' server dan dijawab 80 milidetik kemudian. \\
\term{Jitter} & Seberapa \textbf{tidak stabil} \textit{delay}-nya. Kalau \textit{delay} kadang 50 ms, kadang 400 ms, \textit{jitter}-nya besar. Mirip angkot yang kadang datang 5 menit, kadang 40 menit. \\
\term{Packet loss} & Persentase \textbf{potongan data yang hilang} di jalan. Di penelitian ini diganti dengan \textbf{tingkat gagal}, yaitu persentase pengiriman yang gagal, karena aplikasi tidak bisa melihat paket satu per satu. \\
\term{Throughput} & \textbf{Kecepatan nyata} pengiriman data, misalnya berapa kbps saat foto absen diunggah. Makin besar makin bagus. \\
\term{TIPHON} & \textbf{Standar penilaian} kualitas jaringan dari lembaga Eropa (ETSI). Angka hasil ukur dicocokkan dengan tabelnya untuk mendapat kategori ``sangat bagus'', ``bagus'', ``sedang'', atau ``buruk''. \\
\term{Wi-Fi / seluler} & Wi-Fi = internet dari router kantor. Seluler = internet dari kartu SIM (4G). \\
\term{Blank spot} & Area yang \textbf{tidak ada sinyal sama sekali}. \\
\multicolumn{2}{l}{\textbf{B. Istilah pengukuran}}\\
\term{Crowdsourcing} & Mengumpulkan data dari \textbf{banyak orang sekaligus}. Di sini ``orang banyak'' adalah HP karyawan yang dipakai absen. Mirip Google Maps yang tahu jalan macet dari HP banyak pengguna, atau OpenSignal/Ookla yang memetakan sinyal dari HP pengguna. Perannya: \textbf{teknik pengumpulan data}. \\
\term{Opportunistic sensing} & Jenis \textit{crowdsourcing} yang datanya diambil \textbf{otomatis tanpa tindakan pengguna}. Kebalikannya \textit{participatory sensing}, yaitu pengguna sengaja menekan tombol untuk mengukur. Penelitian ini memakai \textit{opportunistic sensing}. \\
\term{Telemetri} & Dari bahasa Yunani \textit{tele} (jauh) + \textit{metron} (ukur): pengukuran \textbf{otomatis} dari jarak jauh, lalu hasilnya dikirim ke server. Contoh: mobil F1 mengirim data mesin ke tim di pit. Perannya: \textbf{instrumen penelitian}. \\
\term{Pasif} & Mengukur dari kegiatan yang \textbf{memang sudah terjadi} (aplikasi menghubungi server), bukan menjalankan tes khusus seperti Speedtest. \\
\term{Aktif} & Mengirim permintaan khusus untuk mengukur, seperti \textit{probe}. Penelitian ini hibrida: sebagian besar pasif, ditambah \textit{probe} aktif yang sangat ringan. \\
\term{RTT (Round-Trip Time)} & Waktu pergi-pulang data dari HP ke server dan kembali. Di penelitian ini diukur di tingkat aplikasi (HTTP). \\
\term{HTTPS / TCP / TLS} & HTTPS = protokol aplikasi untuk berkomunikasi dengan server secara terenkripsi. Berjalan di atas TCP (pengiriman data yang andal) dan TLS (enkripsi). \\
\term{Probe} & ``Ketukan'' kecil ke server (3 kali, paling sering tiap 15 menit) untuk mengukur \textit{delay} dan \textit{jitter}. Ukurannya sangat kecil. \\
\term{Anonim (hash)} & ID karyawan diubah menjadi kode acak yang tidak bisa dikembalikan ke nama. Tidak ada nama, GPS, atau foto yang dicatat. \\
\term{Saklar server} & Pengaturan di server untuk \textbf{menyalakan/mematikan} telemetri di semua HP tanpa update aplikasi. \\
\term{PCAPdroid} & Aplikasi Android untuk \textbf{merekam lalu lintas data}. Dipakai untuk mengecek apakah angka telemetri akurat. \\
\term{Validasi instrumen} & Membuktikan \textbf{alat ukur buatan sendiri} memberi angka yang benar, dengan membandingkannya dengan alat yang sudah diakui. \\
\term{Server uji (staging)} & Server cadangan khusus untuk mencoba-coba, \textbf{terpisah} dari server yang dipakai karyawan. \\
\multicolumn{2}{l}{\textbf{C. Istilah K-Means}}\\
\term{Clustering (klasterisasi)} & Mengelompokkan data yang \textbf{mirip ke dalam satu kelompok} secara otomatis, tanpa diberi label dulu. \\
\term{K-Means} & Algoritma \textit{clustering} paling populer. ``K'' adalah jumlah kelompok. Cara kerjanya: tentukan K titik pusat, masukkan setiap data ke pusat terdekat, geser pusat ke rata-rata anggotanya, ulangi sampai stabil. \\
\term{Klaster} & Satu kelompok hasil K-Means, misalnya ``jaringan buruk''. \\
\term{Pusat klaster (centroid)} & Nilai rata-rata satu kelompok, dipakai untuk memberi nama kelompok (misalnya \textit{delay} tinggi = buruk). \\
\term{Metode elbow} & Cara memilih jumlah kelompok terbaik: buat grafik untuk K = 2 sampai 6, cari titik ``siku'' tempat penurunan grafik mulai melandai. \\
\term{Silhouette} & Nilai -1 sampai 1 yang menunjukkan \textbf{seberapa rapi} pengelompokan. Mendekati 1 berarti kelompok jelas terpisah. \\
\term{Normalisasi} & Menyamakan skala angka (misalnya milidetik dan kbps) sebelum K-Means, agar satu metrik tidak mendominasi. \\
\multicolumn{2}{l}{\textbf{D. Istilah metodologi}}\\
\term{R\&D} & Metode penelitian berupa \textbf{proses membuat lalu menguji} produk. Produknya di sini modul telemetri. \\
\term{Borg dan Gall} & Nama \textbf{model R\&D} yang paling sering dipakai. Aslinya 10 tahap (kumpulkan informasi, rencana, buat produk awal, uji coba, revisi, \ldots, sebarluaskan). Di penelitian ini diringkas menjadi 5 tahap: analisis, rancang, implementasi, validasi, pengumpulan \& analisis. \\
\term{Populasi / sampel} & Populasi = semua komunikasi aplikasi ke server. Sampel = yang tercatat selama 4 minggu pengumpulan. \\
\term{GAP penelitian} & Celah yang belum dijawab penelitian sebelumnya. \\
\term{Limitation} & Batasan penelitian, yaitu hal yang tidak ikut diukur atau dikendalikan. \\
\bottomrule
\end{longtable}

% ------------------------------------------------------------------
\section{Cara Pengujian (Dari Meja + Dari Karyawan)}
\begin{enumerate}
\item \textbf{Di meja (validasi):} pasang aplikasi versi uji yang terhubung ke server uji. Buat sendiri kondisi Wi-Fi, 4G, mode pesawat, dan putus-sambung. Bandingkan angka telemetri dengan PCAPdroid.
\item \textbf{Dari karyawan (data utama):} rilis update aplikasi dengan telemetri dalam keadaan \textbf{mati}. Setelah aman, nyalakan dari server untuk beberapa orang dulu, lalu semua. Data terkumpul sendiri selama 4 minggu.
\item \textbf{Analisis di laptop:} unduh data dari server sebagai CSV, lalu olah dengan Python (pandas, scikit-learn) untuk tabel TIPHON dan K-Means.
\item \textbf{Interaksi dengan orang lain:} izin tertulis perusahaan dan satu wawancara singkat dengan HRD (bisa lewat chat).
\end{enumerate}
\textbf{Kenapa tidak mengganggu karyawan:} telemetri hanya mengamati, tidak mengubah cara absen. Kalau ada kesalahan, modul gagal diam-diam tanpa memengaruhi absen, dan bisa dimatikan dari server dalam hitungan detik.

% ------------------------------------------------------------------
\section{Kemungkinan Pertanyaan Dosen}
\subsection{``Apa topik keilmuan dan metodenya?''}
Topiknya jaringan (QoS), tekniknya K-Means, dan metodenya R\&D: membuat modul telemetri, memvalidasinya, lalu memakainya untuk mengumpulkan data.

\subsection{``Judulnya `Analisis', kenapa metodenya R\&D?''}
Penelitian ini mengembangkan modul telemetri (produk R\&D) yang kemudian dipakai sebagai instrumen untuk menganalisis QoS jaringan. Kata ``Analisis'' merujuk pada tahap akhir penelitian, yaitu analisis data hasil uji coba pemakaian produk.

\subsection{``Crowdsourcing itu metodenya?''}
Bukan. \textit{Crowdsourcing} adalah \textbf{teknik pengumpulan data}. Metode penelitiannya R\&D, instrumennya modul telemetri, dan metode analisisnya K-Means (lihat bagian 2).

\subsection{``Protokol apa yang kamu ukur?''}
Komunikasi HTTPS (HTTP di atas TCP/TLS) antara aplikasi dan server, lewat Wi-Fi (IEEE 802.11) atau seluler. \textit{Delay} yang diukur adalah RTT tingkat aplikasi, lalu divalidasi dengan RTT tingkat TCP dari PCAPdroid (lihat bagian 3).

\subsection{``Apa bedanya dengan penelitian QoS lain?''}
Penelitian lain mengukur manual di satu tempat dan satu waktu. Penelitian ini mengukur otomatis dari banyak HP, banyak lokasi, dan seluruh jam kerja selama 4 minggu.

\subsection{``Apakah boleh mengambil data dari HP karyawan?''}
Boleh dengan izin tertulis perusahaan. Data dibuat anonim (tidak ada nama, GPS, atau foto), sesuai prinsip UU Pelindungan Data Pribadi No. 27/2022, dan ada pemberitahuan di menu Tentang Aplikasi.

\subsection{``Bagaimana memastikan angka telemetri benar?''}
Divalidasi di meja: angka telemetri dibandingkan dengan PCAPdroid pada kondisi jaringan yang sama. Target selisih paling banyak 20\%.

\subsection{``Kenapa packet loss diganti tingkat gagal?''}
Aplikasi tidak bisa melihat paket satu per satu tanpa akses khusus. Persentase pengiriman yang gagal adalah pendekatan yang bisa diukur dari dalam aplikasi, dan ini dicatat sebagai batasan.

\subsection{``Kenapa K-Means, bukan cukup tabel TIPHON?''}
TIPHON menilai satu metrik satu per satu. K-Means melihat semua metrik \textbf{sekaligus} dan menemukan pola kelompok lokasi/jam secara otomatis. Hasil keduanya dibandingkan.

\subsection{``Kenapa referensi TIPHON tahun 1999?''}
TIPHON adalah \textbf{dokumen standar}, bukan jurnal penelitian, dan masih menjadi acuan di jurnal QoS terbaru (misalnya Ramayanti 2024, Fiqria 2024, Basmar 2026).

% ------------------------------------------------------------------
\section{Daftar Referensi dan Isinya}
\begin{longtable}{KV}
\toprule
\kepala{Referensi}{Dipakai untuk}
\endhead
\multicolumn{2}{l}{\textbf{A. QoS dan TIPHON}}\\
\term{ETSI TIPHON (1999)} & Standar kategori QoS. \\
\term{Ramayanti (2024)} & Contoh pengukuran QoS jaringan nirkabel dengan TIPHON. \\
\term{Fiqria \& Purnama (2024)} & Contoh QoS memakai Wireshark (pengukuran manual satu titik, mendukung GAP). \\
\term{Rosid dkk. (2023)} & Rumus dan langkah menghitung parameter QoS. \\
\term{Basmar \& Yusuf (2026)} & Bukti kualitas jaringan turun pada jam sibuk (alasan mengukur sepanjang hari). \\
\multicolumn{2}{l}{\textbf{B. Crowdsourcing dan pengukuran pasif}}\\
\term{Wamser dkk. (2021)} & Pengukuran jaringan \textit{crowdsourcing} dan berapa banyak sampel yang dibutuhkan agar valid. \\
\term{Wu dkk. (2017)} & MopEye: mengukur kualitas jaringan secara pasif dari lalu lintas aplikasi. \\
\multicolumn{2}{l}{\textbf{C. K-Means}}\\
\term{Irwansyah dkk. (2025)} & K-Means untuk mengelompokkan kualitas internet seluler (data Ookla). \\
\term{Mubarak dkk. (2023)} & K-Means untuk mengelompokkan pemakaian \textit{bandwidth}. \\
\multicolumn{2}{l}{\textbf{D. Konteks aplikasi dan metode}}\\
\term{Indah \& Wati (2025)} & Aplikasi absensi yang butuh mode offline karena jaringan. \\
\term{Rokhman dkk. (2026)} & Aplikasi presensi/personalia Flutter. \\
\term{Jamaludin dkk. (2024)} & Definisi metode R\&D dan 10 tahap model Borg dan Gall, diterapkan pada pengembangan aplikasi Android (Teknik Informatika, Universitas Raharja). \\
\bottomrule
\end{longtable}

\end{document}
